\documentclass{jsarticle}

\begin{document}

-----------------------
畑中さん
-----------------------
\begin{enumerate}
    \item  論文で具体的に扱っているSAWの周波数は4$\sim$5GHzなので、その場合の波長は$\lambda_a=$1umかまたは数100nm程度かと思われます。一方で、論文中の計算で改定しているYIG板の膜厚はL$=$500umとのことですので、膜厚はSAW波長より十分に大きい構造となっています($L\gg\lambda_a$)。この場合では、YIGの表面と裏面でのモードが結合して対称・反対称Lambモードを形成するほど強い結合は得られず、むしろ、ほぼ両面で独立したSAWレイリー波に近い波が生じていると考えた方が自然なのではと感じました。もちろん、$L\gg\lambda_a$であっても、SAWのような半無限体を仮定した固有モードでなく、有限膜厚における固有モードなので、Lamb波の一種だと思いますが、その特性はSAWにほぼ漸近していると思われます。一般的にLamb波を考える上では$L\sim\lambda_a$が妥当な構造ではと感じました。L=数umで計算した場合では、Fig2と似たような結果が得られないでしょうか?もしそうならその結果を見せた方が一般的なLamb波を扱う構造に近いですし、おそらく、モードオーバーラップも大きくなるので、より鮮明なピークの反交差が現れるのではと推察します。
    \item P. 4の左側の第4パラグラフで、”… BVWs appear, plotted as blue and red curves …”と記載されていますが、Fig. 2(a)のBVWの分散曲線の青と赤が見えにくいです。黒点をもう少し小さくして、赤・青線の視認性を上げることは可能でしょうか。
    \item Fig. 2aで、BVWの基底モードである最低周波数は、膜厚方向に一様なモード分布形状で、symmetricな対称モードになるかおもいました(勘違いでしたら申し訳ございません)。ただ、計算では赤線のasymmetricな反対称モードになっているようです。こちらは膜厚方向でどのようなモード形状を持っていますか?
\end{enumerate}
-----------------------
畑中さん
-----------------------
\begin{enumerate}
    \item この理論構築の新規部分をもう少し主張した方が良いように感じます。Kalnikos-Slavinが与えたダイポール相互作用に関するグリーン関数を変位場で摂動展開し、変位場とマグノンに対する結合方程式を厚み方向に対して等価に扱った点が新規という理解でしょうか?あるいはそれを数値的に解いた部分でしょうか?既存理論枠組みで議論されてきた論文を引用しながら、今回の理論構築で新たに予言されること、新たな適用範囲などをイントロ部などで整理したほうが良いかと思いました。
    \item 数値結果で対称性から結合の様子を議論するのはよくわかるのですが、分裂幅several MHzの正当性を議論する必要は無いでしょうか?マグノメカニクスの潮流的にこのあたり実験屋さんが一番気になるところかと思います。数GHzの共振器でQ値1000くらいを仮定するとラフには強結合領域ですかね。
    \item 45度の結果で中央下部のインセットですが、マグノン側は2本の線がほとんど重なって交差するように見受けられますが、これはそのうちの1本との交差を示しているということでしょうか。あとFig.2の横軸のkLのラベルが全体的に右側にずれているように見受けられますが中央に表示した方が良いかと思います(kL$=$1の部分に書いている?)。
    \item ひょっとすると今回の論文から敢えて外されているのかもしれませんが、以前の打ち合わせでコメントさせていただいた通り、既存理論の磁気弾性定数と今回の理論で算出されるものの対応関係は調べたほうが良いかと思いますがいかがでしょうか。
    \item あまり本質的ではありませんが、L$=$500umでThin filmというのは少し違和感がありました。
    \item 完全に数式を追い切れておりませんで恐縮ですが、磁気弾性結合によるフォースfに近似なしでikが出てくるのはLamb波だからという理解で正しいでしょうか。通常のGivenな磁気弾性結合$\epsilon_{ij} x_i x_j$から出すときはフォノンの分布に応じてGeneralな微分項(例えば$\partial u_x/\partial z$など)が出てきて、近似的にikに比例する項だけ残す営みかと思いましたので質問させていただきました。
\end{enumerate}


\end{document}